\section{Filas}

Uma fila, assim como a pilha, é uma estrutura de dados que implementa uma política para
a ordem de inserções e remoções no conjunto.
A política implementada para a fila é a \textbf{FIFO}: o primeiro elemento a entrar
é também o primeiro a sair, assim como filas de atendimento.
O ato de inserir numa fila é chamado de \textit{enqueue}, enquanto o de remover de
\textit{dequeue}.
Similarmente às pilhas, elas podem ser implementadas a partir de vetores ou listas
encadeadas.

A forma mais intuitiva de criar uma fila com um vetor é inserir os novos membros
num extremo, e remover a partir do outro.
Entretanto, supondo que as inserções são feitas no início e as remoções no fim do vetor,
toda inserção implicaria na movimentação de todos os elementos já presentes.
O mesmo valeria caso as remoções fossem no início.

Uma abordagem possível para realizar as operações de \textit{enqueue} e \textit{dequeue}
em tempo constante é fazendo com que a fila seja circular, de forma a utilizar dois
índices especiais: $f$ e $r$, que indicam, respectivamente, o primeiro e o último elemento
na fila.
Quando a fila está inicialmente vazia, temos $f=0$ e $r = n-1$, em que $n$ é a capacidade
da fila.
Ao adicionar um novo membro, executamos $r \gets (r+1)\bmod n$ e inserimos o elemento
na posição $r$ do vetor.
A operação de \textit{dequeue} é realizada retornando o membro apontado por $f$ e depois
fazendo $f \gets (f + 1) \bmod n$.

\begin{algorithm}
\SetAlgoLined
\Fn{\Enqueue{$Q$, $x$}}{
    \Se{$Q.\NumElem = Q.\Len$}{
        \textbf{erro} ``overflow''\;
    }
    $Q.r \gets (Q.r + 1) \bmod Q.\Len$\;
    $Q[Q.r] \gets *x$\;
    $Q.\NumElem \gets Q.\NumElem + 1$\;
}

\Fn{\Dequeue{$Q$}}{
    \Se{$Q.\NumElem = 0$}{
        \textbf{erro} ``underflow''\;
    }
    $x \gets Q[Q.f]$\;
    $Q.f \gets (Q.f + 1) \bmod Q.\Len$\;
    $Q.\NumElem \gets Q.\NumElem - 1$\;
    \Retorna{$x$}\;
}
\end{algorithm}

A abordagem com listas encadeadas é substancialmente mais simples e intuitiva.
A adição dum elemento novo é feita sempre ao final da própria lista, enquanto
a remoção é feita no início.
A razão é a mesma explicada para as pilhas: não é possível realizar a remoção
do último elemento da lista sem buscar o elemento que o antecede, fazendo com que
a operação não seja realizada em tempo constante.

\begin{algorithm}
\SetAlgoLined
\Fn{\Enqueue{$Q$, $x$}}{
    \InsereLista{$Q$, $Q.\Tail$, $x$}\;
}

\Fn{\Dequeue{$Q$}}{
    \Se{$Q.\Head.\Prox = \nil$}{
        \textbf{erro} ``underflow''\;
    }
    $x \gets *(Q.\Head.\Prox)$\;
    \RemoveLista{$Q$, $Q.\Head$}\;
    \Retorna{$x$}\;
}
\end{algorithm}
